% ============================================================================
% SECTION 3: THE S-BOX STATISTICAL MODEL
% ============================================================================
\section{The S-box Statistical Model}\label{sec:manifold}

\subsection{Key-induced boundary distributions}

Fix a bijective S-box $S:\B{s}\to\B{s}$ and let $A$ be uniformly distributed
over $\B{s}$.  Under key hypothesis $k\in\B{s}$, the observed input-output pair
is $(A,S(A\oplus k))$.  This induces the following distribution on
$\mathcal X=\B{s}\times\B{s}$.

\begin{definition}[Key-induced distribution]\label{def:key-dist}
For $k\in\B{s}$, define
\begin{equation}\label{eq:pk}
  P_k(a,b)=2^{-s}\mathbf 1[b=S(a\oplus k)].
\end{equation}
The support of $P_k$ is the graph
$G_k=\{(a,S(a\oplus k)):a\in\B{s}\}$.
\end{definition}

Each $P_k$ is a boundary point of the simplex $\Delta_{2^{2s}-1}$, because it
assigns zero mass to $2^{2s}-2^s$ outcomes.  This support deficiency is central:
ordinary Fisher--Rao tensors of the full simplex are defined on the interior,
not directly at $P_k$.  We therefore treat $P_k$ as a singular endpoint of an
interior regularization when differential-geometric operations are needed.

\subsection{Walsh representation and key translation}

\begin{definition}[Walsh coefficient of $P_k$]\label{def:wht-pk}
For $(\alpha,\beta)\in\B{s}\times\B{s}$, set
\begin{equation}\label{eq:wht-pk}
  \widehat P_k(\alpha,\beta)
  =\sum_{(a,b)\in\mathcal X}\chi_{\alpha\beta}(a,b)P_k(a,b)
  =2^{-s}\sum_{a\in\B{s}}(-1)^{\alpha\cdot a\oplus\beta\cdot S(a\oplus k)}.
\end{equation}
\end{definition}

\begin{lemma}[Walsh key-translation rule]\label{thm:phase-shift}
For every $k\in\B{s}$ and every $(\alpha,\beta)\in\B{s}\times\B{s}$,
\begin{equation}\label{eq:phase-shift}
  \widehat P_k(\alpha,\beta)
  =(-1)^{\alpha\cdot k}\widehat P_0(\alpha,\beta)
  =2^{-s}(-1)^{\alpha\cdot k}\widehat S(\alpha,\beta).
\end{equation}
\end{lemma}

\begin{proof}
Using $a'=a\oplus k$,
\[
\widehat P_k(\alpha,\beta)
=2^{-s}\sum_{a'}(-1)^{\alpha\cdot(a'\oplus k)\oplus\beta\cdot S(a')}
=2^{-s}(-1)^{\alpha\cdot k}\sum_{a'}(-1)^{\alpha\cdot a'\oplus\beta\cdot S(a')},
\]
which is exactly the claimed identity.
\end{proof}

\begin{remark}[Walsh orbit]
The key does not change the amplitude $|\widehat S(\alpha,\beta)|$; it only
changes the sign according to the character $\alpha\mapsto(-1)^{\alpha\cdot k}$.
Thus $\{P_k:k\in\B{s}\}$ is best described as a finite Walsh-phase orbit, not as
a smooth manifold.  The rule is a structural translation identity; the
information-geometric content in the sequel comes from the covariance limit,
active subspace, and belief-update geometry built on top of this identity.
\end{remark}

\subsection{The regular Walsh exponential family}

\begin{definition}[Walsh exponential family]\label{def:walsh-exp}
Let $T_{\alpha\beta}=\chi_{\alpha\beta}$ for
$(\alpha,\beta)\neq(0,0)$.  The Walsh exponential family on $\mathcal X$ is
\begin{equation}\label{eq:walsh-exp}
  p_\theta(a,b)=
  \exp\!\left(\sum_{(\alpha,\beta)\neq(0,0)}
  \theta_{\alpha\beta}T_{\alpha\beta}(a,b)-\psi(\theta)\right),
  \qquad \theta\in\R^{2^{2s}-1}.
\end{equation}
\end{definition}

Because the nonconstant Walsh characters together with the constant character
form an orthogonal basis of all real functions on $\mathcal X$, this is simply
the full regular exponential family over $\mathcal X$ written in Walsh
coordinates.  Its image is the interior of the probability simplex.  The
singular distributions $P_k$ belong to the closure, not to the regular
parameter space.

\begin{proposition}[Boundary expectation coordinates]\label{prop:eta-coords-boundary}
For each $k$ and $(\alpha,\beta)\neq(0,0)$,
\begin{equation}\label{eq:eta-params}
  \eta^{(k)}_{\alpha\beta}
  :=\mathbb E_{P_k}[T_{\alpha\beta}]
  =(-1)^{\alpha\cdot k}2^{-s}\widehat S(\alpha,\beta).
\end{equation}
The corresponding natural parameters are not finite; they are obtained only as
limits of full-support approximations approaching $P_k$.
\end{proposition}

\begin{proof}
The first equality is the definition of expectation coordinates and the second
is Lemma~\ref{thm:phase-shift}.  Since $P_k$ has zeros on $\mathcal X$, no
finite vector $\theta$ in~\eqref{eq:walsh-exp} can represent it.
\end{proof}

\subsection{Interior regularization and boundary covariance form}\label{subsec:fim}

Let $U$ denote the uniform distribution on $\mathcal X$.  For
$\varepsilon\in(0,1)$ define
\begin{equation}\label{eq:eps-reg}
  P_k^{\varepsilon}=(1-\varepsilon)P_k+\varepsilon U.
\end{equation}
Then $P_k^{\varepsilon}$ has full support and belongs to the regular simplex
interior.  Fisher--Rao calculations are unambiguous at $P_k^{\varepsilon}$.
The finite covariance limit as $\varepsilon\downarrow0$ is the following.

\begin{definition}[Boundary Fisher covariance form]\label{def:boundary-fisher}
The boundary covariance form associated with key $k$ is
\begin{equation}\label{eq:boundary-fisher}
  \mathcal I^{\partial}_{ij}(k)=\operatorname{Cov}_{P_k}(T_i,T_j),
  \qquad i,j\in\B{s}\times\B{s}\setminus\{(0,0)\}.
\end{equation}
It is positive semidefinite and generally singular.  We do not call it a
non-degenerate Fisher metric at $P_k$.
\end{definition}

\begin{proposition}[Closed form of the boundary covariance]\label{thm:fim}
For $i=(\alpha,\beta)$ and $j=(\alpha',\beta')$,
\begin{align}
  \mathcal I^{\partial}_{ij}(k)
  &=\operatorname{Cov}_{P_k}(T_{\alpha\beta},T_{\alpha'\beta'})\notag\\
  &=(-1)^{(\alpha\oplus\alpha')\cdot k}
  \left(
  2^{-s}\widehat S(\alpha\oplus\alpha',\beta\oplus\beta')
  -2^{-2s}\widehat S(\alpha,\beta)\widehat S(\alpha',\beta')
  \right).\label{eq:fim-full}
\end{align}
In particular,
\begin{equation}\label{eq:fim-diag}
  \mathcal I^{\partial}_{(\alpha\beta),(\alpha\beta)}(k)
  =1-\left(2^{-s}\widehat S(\alpha,\beta)\right)^2.
\end{equation}
\end{proposition}

\begin{proof}
Since $T_iT_j=T_{i\oplus j}$,
\[
 \mathbb E_{P_k}[T_iT_j]
 =\widehat P_k(\alpha\oplus\alpha',\beta\oplus\beta')
 =(-1)^{(\alpha\oplus\alpha')\cdot k}2^{-s}
 \widehat S(\alpha\oplus\alpha',\beta\oplus\beta').
\]
Subtracting
$\mathbb E_{P_k}[T_i]\mathbb E_{P_k}[T_j]$ and applying
Lemma~\ref{thm:phase-shift} gives~\eqref{eq:fim-full}.  The diagonal follows
from $T_i^2=1$.
\end{proof}

\begin{lemma}[Key covariance isometry]\label{lem:key-isometry}
Let $D_k$ be the diagonal sign matrix
$(D_k)_{(\alpha\beta),(\alpha\beta)}=(-1)^{\alpha\cdot k}$.  Then
\begin{equation}\label{eq:fim-factor}
  \mathcal I^{\partial}(k)=D_k\mathcal I^{\partial}(0)D_k.
\end{equation}
Thus every spectral invariant of the boundary covariance form is independent of
the true key.
\end{lemma}

\begin{proof}
Equation~\eqref{eq:fim-full} factors exactly into the sign contribution
$(-1)^{\alpha\cdot k}(-1)^{\alpha'\cdot k}$ times the $k=0$ covariance entry.
\end{proof}

\begin{remark}[What is metric and what is not]
For $\varepsilon>0$, the Fisher--Rao metric at $P_k^\varepsilon$ is an ordinary
Riemannian metric.  At $\varepsilon=0$, only the finite covariance form
\eqref{eq:boundary-fisher} is used.  Distances or curvatures involving the
exact $P_k$ are therefore either (i) computed on the belief simplex,
(ii) computed for $P_k^\varepsilon$ and then studied as limits, or
(iii) explicitly described as first-order/local covariance approximations.
\end{remark}

\subsection{First-order spectral separation of key hypotheses}

The exact distributions $P_k$ and $P_{k'}$ have disjoint supports when
$k\neq k'$, so their KL divergence is infinite.  Nevertheless, the Walsh
coordinate displacement has a simple closed form:
\begin{equation}\label{eq:eta-displacement}
  \eta^{(k')}_{\alpha\beta}-\eta^{(k)}_{\alpha\beta}
  =\left((-1)^{\alpha\cdot k'}-(-1)^{\alpha\cdot k}\right)2^{-s}\widehat S(\alpha,\beta).
\end{equation}
Consequently, only Walsh modes satisfying $\alpha\cdot(k\oplus k')=1$ contribute
to the spectral separation.  A diagonal local approximation based on the
boundary covariance form is
\begin{equation}\label{eq:fr-local}
  d_{\mathrm{loc}}^2(k,k')
  =\sum_{\substack{(\alpha,\beta)\neq(0,0):\\
  \alpha\cdot(k\oplus k')=1}}
  \frac{4\,2^{-2s}\widehat S(\alpha,\beta)^2}
  {1-2^{-2s}\widehat S(\alpha,\beta)^2},
\end{equation}
provided the denominator is nonzero.  This quantity is a local spectral proxy,
not the exact Fisher--Rao geodesic distance between the singular distributions.
